\section{Introduction}\label{sec:introduction}

Maritime transportation is the driving engine of the global economy, representing 80 \% of the world trade by value~\citep{Unctad2023}. It is the most energy-efficient mode of transportation compared to air, rail, or road, thanks to its carrying capacity and economical fuel type. An average-size cargo shipment is comparable to 2000 trucks, 2500 airplanes, or 225 trains as mentioned by Pena et al.~\citep{penaReviewApplicationsMachine2020}.

Much of the emphasis recently has been as highlighted by Ji et al.,Pavlenko et al.,Stolz et al.~\citep{jiDatadrivenStudyIMO2020, pavlenkoClimateImplicationsUsing, stolzTechnoeconomicAnalysisRenewable2022} on the usage of clean energy such as \ac{vlsfo}, \ac{mgo} and \ac{lng}. 
However, according to the 4th \ac{ghg} study \ac{imo}~\citep{imoFourthIMOGHG2020}, 
most vessels today still uses \ac{hfo}, which is roughly 30\% cheaper compared to other fuel types since \ac{hfo} is the residual product of the oil refinement process for the extracption of gasoline and diesel so-called distillates for use in planes, trucks and cars which leaves it as waste product and this makes it cost-effective~\citep{clearseasMarineFuelsWhat2020}. 

This fuel is used in the ship in four main areas: main engine propulsion, energy saving for main engine, auxiliary propulsion, and energy saving for auxiliary engine as mentioned by Fan et al.~\citep{fanReviewShipFuel}. However, although shipping is a very efficient mode of transport, it has a significant footprint on the climate; this is multiplied as the volume of international trades that rely on shipping is rising. Roughly 3\% of all global \ac{ghg} emissions are caused by maritime transportation~\citep{Unctad2023}. To put this into perspective, that is 1 Billion tons of \ac{ghg} emitted annually to the atmosphere, which is equivalent to the emissions of whole countries such as Mozambique, for example, with 28 million habitats Pena et al.~\citep{penaReviewApplicationsMachine2020}.

To reduce this, \ac{imo} have proposed regulations that pressure the industry to drive change. One of the leavers in the context of \ac{ghg} is the ambition to lower emissions by half by the year 2050 compared to levels of 2008 which was agreed upon in the Paris Agreement on climate change in~\citep{unParisAgreement2015} with the overall target of lowering global temperature to less than 2° globally Pena et al.~\citep{penaReviewApplicationsMachine2020}.

This shared goal puts an urgent need to optimize voyages to meet \ac{imo} goals. All optimization efforts are classified into two categories: design optimization, for example, wind-assisted propellers or hull-fouling, and operational optimization, such as just-in-time arrival, slow-steaming, and green routing.

During operation there exists a strong connection between \ac{eeoi}, \ac{cii} and \ac{fc} Ailong et al.,Zhang et al.,Yapeng et al.~\citep{FAN2021120266, zhang2017multi, HE2021109733} as cited by Fan et al.~\citep{fanReviewShipFuel}. Therefore, the reduction of fuel can result in significant operational cost savings. Hence, optimal \ac{fc} is crucial for realizing this. Also, as pointed out by the latest review paper Ran et al.~\citep{YAN2021102489} as cited by Li et al.~\citep{liDataFusionMachine2022}, the basis of all operational measures for ship bunker fuel savings and emission mitigation is the quantitatively modeling the relationship between \ac{fc} rate and its determinants, including sailing speed, draft/displacement, trim, weather conditions, and sea conditions, but it is not a trivial work.

The accurate predictions of \ac{fc} come down to selecting a suitable model and data that captures all the influences on the bunker fuel. A recent review paper highlighted that Fan et al.~\citep{fanReviewShipFuel} choice of models could be classified into two types: (a) \ac{wbm}, sometimes referred to as physics-based models, (b) \ac{bbm}, sometimes referred to as \ac{ml} models.

White-box models are methods that are based on ship and engine-propeller performance models; see review Ran et al.~\citep{YAN2021102489} and Fan et al.~\citep{fanReviewShipFuel}.
They are based on accurately estimating total ship resistance Zhang et al.~\citep{zhangDeepLearningMethod2024}. The resistance that the ship experiences during sailing could be classified into two types: (a) Calm water resistance and (b) In-service resistance, sometimes referred to as added resistance Fan et al.~\citep{fanReviewShipFuel}.
Calm water resistance refers to the hydrodynamic effects the ship encounters when sailing in undisturbed waters, such as viscous or wave-making resistance, which are generally caused by the ship. In contrast, added resistance is based on the environment. It refers to the hydrodynamics effects on a ship from external factors such as winds, waves, other ships, and obstacles Vinayak et al.,Zhihua et al.~\citep{vinayaketal2021, LIU2020107246} as cited by Zhang et al.~\citep{zhangDeepLearningMethod2024}. The total resistance of a ship is computed by adding both of those components, and the most common methods used are the towing tank and \ac{cfd} Fan et al.~\citep{fanReviewShipFuel}. \ac{wbm}, as the name suggests, are transparent in their structure and interpretable in nature. They rely on physical principles, theoretical frameworks, and explicit mathematical formulations to explain the relation between Ship \ac{fc} and its determinants. They are often used in the ship design phase to manufacture efficient ships.
The parameters used in those classes of models are inter-connected and, therefore, susceptible to the environment, which leads to errors, especially in the case of extreme navigation environments Nan et al.~\citep{WEI2021121036} as cited by Fan et al.~\citep{fanReviewShipFuel}. Additionally, the parameters derived in those formulas are not adjustable, which means they cannot be changed during the voyage nor account for degradation over time in the propulsion systems Fan et al.~\citep{fanReviewShipFuel}.
The performance of the \ac{wbm} is strongly affected by various assumptions Ran et al.~\citep{YAN2021102489} as cited by Fan et al.~\citep{fanReviewShipFuel}.
For example, the components and interactions of the ship resistance are often ignored, resulting in poor applicability of \ac{fc} models Haranen et al.~\citep{haranenWhiteGreyBlackBox} as cited by Fan et al.~\citep{fanReviewShipFuel}

The term "white-box" contrasts with "black-box" models, which do not have transparent internal workings and rely heavily on empirical data and statistical methods without necessarily explaining the underlying physical processes. There is a vast potential for using \ac{ml} methods to model Ship \ac{fc} which can solve the challenges that physics-based models Chen et al.,Lang et al.,Shang et al.~\citep{chenetal2023, lang2023data, shangetal2023} as cited by Zhang et al.~\citep{zhangDeepLearningMethod2024}.
This is because \ac{ml} methods could elucidate the inter-relationship between the measured \ac{fc} and its determinant factors based solely on the observed measurements with the ability to take realistic navigational patterns, ship operational status, weather conditions, and engineering systems specifics Zhang et al.~\citep{zhangDeepLearningMethod2024}.~\ac{ml} excels in decoding intricate patterns, making it particularly advantageous for advanced shipping applications, which is the primary focus of this study.

The remainder of this paper is organized as follows: \cref{sec:literature_analysis} The literature analysis, where there is a narrow focus on the current state of research in terms of Ship \ac{fc} prediction using \ac{ml} approaches, is presented, together highlighting the gaps and contribution of this study. \cref{sec:methodological_approach} cover the methodological approach used, and \cref{sec:results} highlights and discuss the results. Finally, a conclusion of the findings is drawn together with an outlook of future research are presented in \cref{sec:conclusion}.